\chapter{String Representations, Unicode, and String Operations}
\label{ch:strings-representation}

A compiler can scan a source file correctly and still underline the wrong character in a diagnostic.
The scanner counts bytes, the editor counts UTF-16 code units, and the display groups several code points into one visible cluster.
A single integer called \code{length} cannot reconcile those coordinate systems.
String processing begins by choosing what the symbols mean and which operations preserve their interpretation.

\section{Define the sequence before counting it}
\label{sec:ch77-sequences}
For an alphabet $\Sigma$, a string $s\in\Sigma^*$ is a finite sequence of symbols.
We write $|s|$ for its length in that alphabet, $\epsilon$ for the empty sequence, and $s[i:j]$ for the half-open slice from $i$ through $j-1$.
Concatenation is written $uv$; a prefix has the form $s[0:j]$, and a suffix has the form $s[i:|s|]$.
The alphabet might consist of bytes, Unicode scalar values, token kinds, or machine opcodes.
Algorithms on sequences transfer between these domains only when their equality and indexing assumptions transfer too.

Consider the text consisting of A followed by U+00E9.
Its UTF-8 bytes are \code{41 C3 A9}; its scalar sequence has length two.
Replacing U+00E9 with U+0065 followed by U+0301 changes the scalar count to three and the UTF-8 byte count to four without changing the canonically equivalent text.
Table~\ref{tab:ch77-units} separates the counts.
The table uses this particular combining sequence; it is not a claim that every scalar is a display cell.

\begin{table}[htbp]
\caption{Equivalent text can occupy different byte and scalar extents.}
\label{tab:ch77-units}
\centering
\begin{tabular}{lrrr}
\toprule
Representation of A plus accented e & UTF-8 bytes & Scalars & Grapheme clusters \\
\midrule
A, U+00E9 & 3 & 2 & 2 \\
A, U+0065, U+0301 & 4 & 3 & 2 \\
\bottomrule
\end{tabular}
\end{table}

\term{Grapheme segmentation} identifies boundaries approximating user-perceived characters under a specified rule set~\cite{unicode-2026-segmentation}.
It does not determine font shaping, terminal display width, or language-specific collation.
An identifier table usually needs exact normalized-key equality, while cursor motion may need grapheme boundaries and sorting may need a locale-sensitive collation service.
Those are different services and should have different APIs.

\section{Choose storage for the update workload}
\label{sec:ch77-storage}
A flat immutable buffer supports compact sequential scans and simple sharing by slices.
Appending by allocating a new flat buffer copies the old prefix; repeated single-character appends can therefore cost quadratic total time.
A growable builder reserves spare capacity and supports amortized constant-time append under geometric growth, with occasional linear relocation.
A rope stores concatenation as a tree, trading pointer chasing for cheaper structural edits when its balancing and chunk-size policies are maintained.

Figure~\ref{fig:ch77-rope} represents two chunks without pretending that their addresses are contiguous.
With a balanced rope of $C$ chunks, locating a byte by stored subtree lengths costs $O(\log C)$ plus work within the selected chunk.
Without balancing, a chain of appends can make that lookup linear in $C$.
Scalar indexing additionally needs decoding or a per-chunk index; byte lengths alone do not provide scalar offsets.

\begin{figure}[htbp]
\centering
\begin{tikzpicture}[node distance=13mm and 15mm]
\node[basic block] (r) {concat: byte extent 7};
\node[basic block,below left=of r] (a) {chunk: 3 bytes};
\node[basic block,below right=of r] (b) {chunk: 4 bytes};
\draw[cfg edge] (r)--(a);\draw[cfg edge] (r)--(b);
\end{tikzpicture}
\caption{A rope stores structural concatenation and must maintain lengths and balance explicitly.}
\label{fig:ch77-rope}
\end{figure}

Compiler source managers often retain immutable input bytes and store token spans as byte offsets.
That permits exact reproduction of spelling and stable references into one source snapshot.
An incremental editor can use a different mutable representation, but an edit must create a new snapshot identity or translate affected spans.
A span from yesterday's buffer is not valid merely because today's buffer has the same total length.

\section{Validate decoding at the boundary}
\label{sec:ch77-decoding}
A UTF-8 decoder classifies the leading byte, checks the required continuation bytes, reconstructs a scalar, and rejects overlong encodings, surrogate values, and values beyond U+10FFFF.
Algorithm~\ref{alg:ch77-decode} makes the rejection policy explicit.
It is a specification-level decoder sketch, not the implementation used by the Python companion, which delegates decoding to Python's strict UTF-8 reader.
For $B$ bytes a full validation scan costs $O(B)$ and can record scalar-start offsets in $O(S)$ storage for $S$ decoded scalars.

\begin{algorithm}[htbp]
\caption{Decode bytes while recording original scalar boundaries.}
\label{alg:ch77-decode}
\begin{pseudocode}
\Require byte buffer with known extent; strict UTF-8 policy
\Ensure scalar sequence and byte-offset map, or first malformed-byte diagnostic
\State $p\gets0$; initialize output and offset map
\While{$p$ is before the end}
\State classify leading byte and required sequence length
\State check extent and continuation-byte shapes before reading them
\State reconstruct value; reject overlong, surrogate, or out-of-range values
\State append scalar and its original byte offset; advance by sequence length
\EndWhile
\State append final byte extent to the boundary map
\end{pseudocode}
\end{algorithm}

For \code{41 C3 A9}, the boundary map is $[0,1,3]$.
Scalar slice $[1,2)$ therefore maps to byte slice $[1,3)$.
The final extent entry lets a half-open slice ending at the string's end use the same conversion rule as an interior slice.
A malformed input policy that inserts replacement characters must instead define which original byte span each replacement covers; silently switching from strict to replacement decoding changes identifier equality.

\section{Normalization is not source preservation}
\label{sec:ch77-normalization}
\term{Unicode normalization} gives canonically or compatibly equivalent sequences a specified normal form~\cite{unicode-2026-normalization}.
NFC performs canonical decomposition followed by canonical composition; NFKC additionally uses compatibility mappings.
Case folding is another transformation, and it can expand a scalar into several scalars.
For example, the companion's caseless key for \code{Straße} is \code{strasse}.
These transformations can change offsets and spelling; they should not overwrite the original source buffer used for diagnostics.

The spellchecker uses the explicit policy NFC, then case fold, then NFC again.
The final normalization accounts for sequences introduced by case folding.
Algorithm~\ref{alg:ch77-key} keeps the display spelling separate from the lookup key.
The companion merges dictionary frequencies when distinct spellings normalize to the same key and returns canonical keys, not a recovered original presentation.
That is a documented interface choice rather than a general typography policy.

\begin{algorithm}[htbp]
\caption{Construct a caseless lookup key without changing the original text.}
\label{alg:ch77-key}
\begin{pseudocode}
\Require validated scalar sequence $s$ and fixed Unicode normalization/case policy
\Ensure canonical lookup key; original spelling remains available separately
\State $n\gets\operatorname{NFC}(s)$
\State $f\gets\operatorname{casefold}(n)$
\State return $\operatorname{NFC}(f)$
\end{pseudocode}
\end{algorithm}

The cost includes the library's normalization work and the expanded output size, not merely the number of input code points.
Canonical ordering of long combining sequences needs particular care before claiming a linear bound for an arbitrary implementation.
The lab places limits both before and after normalization so an expansion cannot bypass the query-length policy.
Production persistence formats should record the normalization policy and Unicode data version; changing either can merge or split keys and require rebuilding an index.

\section{Compose operations with explicit ownership and coordinates}
\label{sec:ch77-operations}
Table~\ref{tab:ch77-operations} lists useful operation contracts.
Substring search reports positions in the sequence it actually searches.
If that sequence is a normalized copy, its offsets cannot automatically be used to highlight the original text.
A transformation-aware source map may need ranges rather than a one-to-one scalar map because composition merges input scalars and case folding can duplicate them.

\begin{table}[htbp]
\caption{String operations need units, equality policies, and lifetime rules.}
\label{tab:ch77-operations}
\centering
\begin{tabular}{ll}
\toprule
Operation & Contract to state \\
\midrule
slice & byte/scalar/cluster indices; owner retention \\
concatenate & copying versus structural sharing \\
compare & binary, normalized, or locale-sensitive equality/order \\
search & alphabet and coordinate system of reported matches \\
intern & canonicalization policy; table lifetime \\
replace & overlap policy; new offset mapping \\
\bottomrule
\end{tabular}
\end{table}

Figure~\ref{fig:ch77-coordinates} keeps the lookup and diagnostic paths apart.
An intern table can store one canonical key and many source occurrences.
Hash equality is only a candidate filter: unless collisions are excluded by construction, compare the actual keys after hash agreement.
A hash-consed identifier also needs a table-lifetime contract; its numeric ID is not stable across arbitrary rebuilds unless deliberately persisted.

\begin{figure}[htbp]
\centering
\begin{tikzpicture}[node distance=13mm and 12mm]
\node[basic block] (source) {immutable source bytes};
\node[basic block,below left=of source] (span) {original diagnostic spans};
\node[basic block,below right=of source] (key) {decoded normalized keys};
\draw[cfg edge] (source)--(span);\draw[cfg edge] (source)--(key);
\end{tikzpicture}
\caption{Canonical lookup need not discard the original coordinates and spelling.}
\label{fig:ch77-coordinates}
\end{figure}

\section{Exercises}
\label{sec:ch77-exercises}
\begin{enumerate}
\item\label{ex:ch77-offsets} Compute the byte boundary map for A, U+0065, U+0301 encoded in UTF-8 and map scalar slice $[1,3)$ back to bytes.
\item\label{ex:ch77-builder} A naive builder copies its complete prefix on each of $n$ one-byte appends.
Compute the total copied prefix bytes and compare with geometric growth.
\item\label{ex:ch77-key} Work the caseless keys for \code{Straße} and \code{STRASSE} and explain why their source spans must remain distinct.
\item\label{ex:ch77-snapshot} Two source snapshots have equal byte lengths but different text before a stored token span.
Explain why length equality does not validate the old span.
\end{enumerate}

\section{Further reading}
\label{sec:ch77-reading}
Crochemore, Hancart, and Lecroq develop strings as algorithmic sequences~\cite{crochemore-2007-strings}.
The Unicode normalization and segmentation annexes specify the text-level distinctions that byte algorithms alone cannot resolve~\cite{unicode-2026-normalization,unicode-2026-segmentation}.
Chapter~\ref{ch:ffi-marshaling} connects explicit string extents to native interfaces.
